\section{Principle I: Truth}

\begin{quotex}
From Chapter II, ``Principles", of Mes Idées Politiques, by \textbf{Charles Maurras}. 
\end{quotex}

Yes sir, yes ma'am, because ``speech is hard", its efficacy will be sweet; because the book is ``bitter to the tongue", it will be invigorating and healing.

The signposts raised up on the road don't show their direction in a sweet and florid style: they flaunt the style of their utility. Clear, direct, insistent, and authoritarian, they don't say: ``if I am mistaken", they don't doubt themselves, they do not excuse themselves from roughly throwing the direction arrows and the mile markers into the eyes of the people who pass. But does the traveler complain about it?  As long as he has the heart of a philosopher, he gives thank to the author of profitable brutalities which he does not even feel tyrannized by.

It is his choice to slow down or step on the gas, to follow to change his direction. The milestone says what it is in clear terms, and what is necessary to take into account. The more that precise facts limit thought and, because of that narrow limit, the more that fantasies of the heart, wishes of the imagination, the needs, the amenities and personal interests will obtain safety and will be able to give themselves a career. An uncertain direction, a fact, whether vague or false, while appearing to flatter the arbitrariness of the walker, will restrict the freedom of his movements, of his rest, they will diminish his real powers, for the risks attached to the consequences of an indifferent or capricious itinerary will be increased by the insufficiency of his knowledge.

It is a great error to think that \emph{contingencies}, as they say, accommodate themselves more easily to a lax and vacillating principle: to the contrary, all indecision of principles complicates the study of the facts, as well as their treatment; uncertainty thus is inserted at the sole point from where a little light could come to them, to the complexities of the earth shadows in the sky will be added.

Truth, a harsh but clear sun, is content to establish from above what is necessary to know and think before acting. It shows the good, it marks out the bad; it distinguishes the proportions following which the one and the other confront each other and mix in the infinite variety of our human events. Once so enlightened, man is far from having resolved the problems of practical life, but he has something to resolve them and if, as happens to him too frequently, he can choose only between evils, he will better discern which will be the least, his effort can be applied to avoiding the worst; that makes perhaps the greatest point of the government of oneself and others.

Not only is truth defended by what it has that is naturally general, elevated, abstract, and foreign to man, but in order to decide to ascend to it, a general impulse [\emph{élan}] of thought beyond the present, a full calculation of the future, is necessary. In order to adhere to this truth that veils everything, in addition an effort of will imposing silence on many instincts is necessary.

Truth (I do not say sincerity, I say the whole truth, the agreement of language and thought with external realities), can still be something besides the highest delight of intelligence: it is the empowered sovereign, it is all powerful force.

Sincerity is not the truth. The most correct intention and the firmest will cannot create what it is not.

Let us not underestimate any virtue, but let us give justice to the obvious virtues. There is no smile, grimace, or chattering of beautiful spirits that can stand against them. The decisions that they lead to are serious, sometimes painful, in the life of the spirit; exterior life doesn't always accommodate them, but the service that they render is such that they prevail over everything.

Truth is of value in itself. But there are bitter truths and sweet truths. There are useful and dangerous truths. There are those that are necessary to reserve only for the wise, and others which are suitable for the nourishment of all.

Some purely oratorical and mystical revolutionaries can believe that, fable or truth, it is always good enough for the people! We believe that the people have no needs less demanding than the elite. Truth is as necessary as bread. The historico-political lie poisons a people as absolutely as potassium cyanide.

A disdain which is not expressed does not take effect. To the contrary, an error and a lie that one does not take the trouble to unmask acquire little by little the authority of the true.

Complete text: \url{http://www.royaliste.org/spip.php?article520}



\flrightit{Posted on 2012-09-30 by Cologero }

\begin{center}* * *\end{center}

\begin{footnotesize}\begin{sffamily}



\texttt{Dominion on 2012-09-30 at 14:45 said: }

"Sincerity is not the truth. The most correct intention and the firmest will cannot create what it is not."

"To the contrary, an error and a lie that one does not take the trouble to unmask acquire little by little the authority of the true."

With those two statements Maurras has unveiled the secret to every victorious lie and legend's success and the key to its power. If a single man or an entire people cannot bring itself to tear the pretty mask off the lies which keep them entranced and captive, they cannot progress. Indeed, they may well fall deeper into the mire.


\hfill

\texttt{Jason-Adam on 2012-10-01 at 17:52 said: }

This text, written in the flowery French style of which Maurras was a master of, is not something easily grasped by those used to modern Joe Friday ``just the facts ma'am" styles of writing. What Maurras is saying here though is just exactly that – before we can figure the world out we need to begin with facts, things we are certain of through proveable means, to avoid being trapped by the lies of the power hungry.


\hfill

\texttt{Attila on 2012-10-21 at 02:24 said: }

I disagree with the characterization of ``flowery French style"— Maurras' writing is highly nuanced, and not reductionist- unlike so much coming from Anglo-Saxon circles.


\end{sffamily}\end{footnotesize}
